\documentclass{article}
\usepackage[left=2cm, right=2cm, top=2cm, bottom=2cm]{geometry}
\usepackage{graphicx}
\usepackage{amsmath}
\usepackage{amssymb}
\usepackage{amsthm}
\usepackage{fancyhdr}
\usepackage{verbatim}
\usepackage{listings}
\usepackage{xcolor}
\usepackage{pdfpages}
\usepackage{cancel}
\usepackage{physics}
\usepackage{esint}
\usepackage{braket}

\lstset{
    language=C++,
    basicstyle=\ttfamily\footnotesize,
    keywordstyle=\color{blue},
    commentstyle=\color{green},
    stringstyle=\color{red},
    numbers=left,
    numberstyle=\tiny\color{gray},
    frame=single,
    breaklines=true,
    captionpos=b,
}


\title{MTH 446: Lecture \# 12}
\author{Cliff Sun}

\newtheorem{theorem}{Theorem}[section]
\newtheorem{lemma}[theorem]{Lemma}
\newtheorem{definition}[theorem]{Definition}
\newtheorem{conjecture}[theorem]{Conjecture}
\newtheorem{proposition}[theorem]{Proposition}
\newtheorem{corollary}[theorem]{Corollary}
\newtheorem{one minute paper}[theorem]{One Minute Paper}

\pagestyle{fancy}
\lhead{\textbf{\thepage}\ \ \nouppercase{\rightmark}}
\chead{MTH 446: Lecture \# 12}
\rhead{Cliff Sun}

\begin{document}

\maketitle

\section*{Lecture Span}
\begin{enumerate}
    \item Complex derivatives
\end{enumerate}

\section*{Complex derivatives}

Note, by the \underline{necessary condition of complex differentiability}, $f'$ DNE for $z$ such that the CR equations are not satisfied. 

By the \underline{sufficient condition of CD}, $f'$ exists for $z$ such that the CR equations are satisfied. 

\section*{CR equations in polar coordinates}

\begin{align}
    u_r &= \frac{1}{r}v_{\theta} \\
    \frac{1}{r}u_{\theta} &= -v_r
\end{align}

\begin{proof}
    Use the chain rule and the fact that $x = r\cos\theta$ and $y = r\sin\theta$ to convert the CR equations in rectangular coordinates to polar coordinates. That is, 
    \begin{equation}
        \frac{du}{d\theta} = u_x x_{\theta} + u_y y_{\theta}
    \end{equation}
\end{proof}

\begin{equation}
    u_x = u_r\cos\theta + \sin\theta/r \cdot u_{\theta}
\end{equation}
\begin{equation}
    u_y = u_r\cos\theta - \sin\theta/r \cdot u_{\theta}
\end{equation}

\subsection*{Example}

If $z_0 \neq 0$, then $f'(z_0)$ exists, then 

\begin{equation}
    f'(z_0) = -\frac{i}{z_0}\left( u_{\theta} +iv_{\theta} \right)
\end{equation}

\begin{proof}
    $f'(z_0) = u_x + iv_x =  u_r\cos\theta - \sin\theta/r \cdot u_{\theta} + i(v_r\cos\theta + -\sin\theta/r \cdot v_{\theta})$. You then obtain 
    \begin{align}
        &= -\frac{i}{r}\left[ u_{\theta}(\cos\theta - i\sin\theta) + iv_{\theta}(\cos\theta - i\sin\theta) \right]\\
        &= -\frac{i}{r}\left[ \left( u_{\theta} + iv_{\theta} \right)e^{-i\theta} \right] = -\frac{i}{re^{i\theta}}\left( u_{\theta} + iv_{\theta} \right) = -\frac{i}{z_0}\left( u_{\theta} + iv_{\theta} \right)
    \end{align}
\end{proof}

\section*{More Operators}

Will not be used in class, but useful to mention for completeness. 

\begin{align}
    \partial_z &= \frac{1}{2}\left( \partial_x - i\partial_y \right) \\
    \partial_{\overline{z}} &= \frac{1}{2}\left( \partial_x + i\partial_y \right) \\
\end{align}

Note that 

\begin{equation}
    \partial_{\overline{z}}f = 0 \iff \text{ CR equations }
\end{equation}

\begin{equation}
    \partial_{z}f = u_x + iv_x
\end{equation}

\section*{Analytical Functions}

\begin{definition}
    Let $D$ be open in $\mathbb{C}$ and not empty. Then, $f$ is \underline{analytic} at $z_0$ if there is $\epsilon >0$ (NOT FOR ALL $\epsilon>0$) such that $D_{\epsilon}(z_0) \subseteq D$ and $f'$ exist for every $z \in D_{\epsilon}(z_0)$. 
\end{definition}

Note, if $f'$ is only differentiable at $z_0$ but DNE for some $z$ close to $z_0$, then $f$ is not analytic at $z_0$.  

\begin{definition}
    $f$ is analytical on $U$ if it is analytic at every point $U$, which is the domain of $f$. Set of all analytical functions is denoted by $H(U)$, $H$ meaning holomorphic. 
\end{definition}

\begin{theorem}
    Let $f$ and $g$ analytic, then 
    \begin{enumerate}
        \item $\alpha f + \beta g \in H(U)$ is analytic
        \item $fg$ is analytic.  
        \item $f/g$ is analytic if $g \neq 0$ at $z_0$. 
    \end{enumerate}
\end{theorem}

\begin{theorem}
    \underline{Uniqueness property of analytic functions}. 
\end{theorem}

\end{document}